% ============================================================
% chunk_A.tex -- pages p000-p019 (front matter + opening of Chapter A)
% Rubinstein, "תבניות שניוניות", Bar-Ilan University, 1970, Part I
% ============================================================

% [front matter: title page]
\begin{center}
% Library markings at top of page (printed library label, not part of
% the book's own title-page design) -- omitted from body, noted here only:
% "אוניברסיטה בר אילן / הספריה למתימטיקה" (library bibliographic label)
% Handwritten accession marks "416308", "2X" -- ignored (handwritten).

{\huge תבניות שניוניות}

\bigskip
\Large
מאת

\bigskip
{\LARGE אלחנן רובינשטיין}

\bigskip
% [Two library rubber stamps appear here in the original scan:
%  one illegible/foreign-script stamp, and one reading
%  "אוניברסיטת בר-אילן / הספריה למתמטיקה ולמדעי המחשב / הוצא מהמלאי" --
%  both omitted from body text as library stamps.]

\bigskip
\large
חלק א'.

\bigskip
תבניות שניוניות ריבועיות ומשוואות פל אלמנטריות.

\vspace{2cm}
\normalsize
כל הזכויות שמורות \hfill ירושלים, שנת תש"ל

% [Another diagonal library stamp appears at the bottom of the page:
%  "אוניברסיטת בר-אילן / הספריה למתמטיקה ולמדעי המחשב" -- omitted.]
\end{center}

\newpage

% [front matter: dedication page]
% Handwritten annotations, a signature, and the date "6.7.50" appear in
% the upper part of this page (ignored -- handwritten marginalia).
% A library stamp "אוניברסיטת בר-אילן / הספריה למתמטיקה ולמדעי המחשב /
% הוצא מהמלאי" also appears (ignored -- rubber stamp).

\begin{center}
\vspace{3cm}
{\large אָסֵף יַרְדֵּן}

\bigskip
מתנת ד"ר דב ירדן ורעיתו רחל

למחלקה למתמטיקה ומדעי המחשב

באוניברסיטת בר-אילן

לזכר אחותו תמה

ובעלה הרב יוסף רוזנטל

ר"מ בישיבת "בית יוסף" בפינסק

וילדם אברהם נתן וילדתם אלתה

שנספו בליטא בשואת תש"א-תש"ב
\end{center}

\newpage

% [front matter: introduction, original page label: I]
\section*{הקדמה}

\noindent 1. תורת התבניות - אם כי אינה אלא פרק מיוחד בתורת המספרים - מהווה את חוט השדרה של כל תורת המספרים.

\medskip
על רוב רובן מהבעיות שנודעו בתורת המספרים אפשר לגשת מנקודת ראות אחת ויחידה - תורת התבניות. לפיכך רבה חשיבותה של תורה זאת ועסקו בזה רוב רובם של המתימטיקאים אשר בהזדמנות זו או אחרת באו במגע עם תורת המספרים.

\medskip
בכל השיטות השונות נהגו המתימטיקאים עד כה להסתמך על רעיון אחד - משותף לכולן. ביסוד הרעיון הונח מושג האי-רציונליות הריבועית והאיקוויולנציה שבין שתי אי-רציונליות ריבועיות.

\medskip
לפיכך נותק כאילו מיד המעבר הטבעי מתורת התבניות הריבועיות אל תורת התבניות השניוניות מסדר גבוה.

\medskip
וזאת היא אחת הסיבות שגרמו להתפתחותה של תורת המספרים ה[?] ראשונים [?], ורק זאת האחרונה היתה עלולה - לפי דעתם של המתימטיקאים - לגשר בין צורת התבניות האלמנטריות לבין החלק העיקרי של התורה, תורת התבניות השניוניות מסדר גבוה.

\medskip
אולם, אם כי תורה זאת היא אחת היצירות הנהדרות שנוצרו אי-פעם במתימטיקה, הרי עבר משבר נפשי עמוק על רבים מבין גדולי המתימטיקאים עד שהחליטו לוותר על הרעיון להמשיך בפוריות ה[?]ישרית [?] הרידונקלית ולהכלילה עד כדי כך, שתצלח למשש מכשיר יעיל בכל שטח של תורת המספרים.

\medskip
גם אעפ"י כן ניכר משבר זה אצל \textbf{דיריכלה}. כי בספרו ניתן למצוא עוד נסיונות לפתרון בעיית אותן התבניות אשר ל[?] בלבד אפשר לנחש שהן שייכות למשפחה אחת גדולה אשר אחד מנציגיה היא המשוואה המפורסמת של תורת התבניות הריבועיות אך ורק על סמך

% [original page: II (inferred; numeral faint/uncaptured in scan)]
רעיפות משטח תורת המספרים האי-רציונליים. אלא, שמוכרח, כנראה, לוותר על רעיון זה, ועל כך מעיד תורת היחידות הידועה של \textbf{דיריכלה}, השייכת כבר לתורת המספרים האלגבראיים, שבה נמצא כבר כר נרחב לפיתוח תורת תבנית היחידה במשלוחה.

\medskip
\noindent 2. בעבודה זאת שמתי לי למטרה לתת את תורת התבניות השניוניות מנקודת ראות מיוחדת.

\medskip
נעשה כאן נסיון מוצלח לבסס את כל התורה של התבניות הריבועיות השניוניות על ידי צמצום התבנית הכללית \underline{לתבנית היחידה}.

\medskip
בדרך זאת מתקבלת שיטה כללית לגבי תבניות שניוניות מסדר כלשהו.

\medskip
אלא שעל-ידי כך עובר מרכז הכובד של כל תורת התבניות לתורת תבנית היחידה, ומכאן החשיבות הנוספת למשוואות פל וכללותיה. כי כמו שיתברר אין כל זכות יתר למשוואה שנהגו עד כה לכנותה בשם "משוואת פל" לגבי תבנית היחידה הכללית בכלל.

\medskip
ומנקודת ראות אחרת שייכות תבניות היחידה לסוג מיוחד של תבניות שעוד מזמנם של גאוס, דיריכלה, ודדקינד נהגו לכנותן בשם תבניות דו-צדדיות.

\medskip
על ידי כך אנו עוברים לשטח חקירה חדש - תורת התבניות הדו-צדדיות.

\medskip
אם כי גאוס, דיריכלה ודדקינד התעמקו אך ורק לתבניות ריבועיות דו-צדדיות נמצאה הגדרה מכלילה למין זה של תבניות, הגדרה המאפשרת להעביר את רוב רובן של המסקנות אשר נקבעו לגבי תבניות ריבועיות מהסוג האמור. יתר על כן: על-ידי העברת ההסתעפות הנ"ל לגבי חבורות שניוניות מסדר גבוה נמצא מספר גדול מקבלות

% [original page: III]
את הצליל האמיתי של משמעותן המלאה.

\medskip
\noindent 3. מתוך התמתחות הנושאים שנגענו בהם מתברר מיד שיש צורך טבעי לתת את תורת התבניות השניוניות בארבעה חלקים, והם:

\medskip
\noindent \underline{חלק א'}: תורת התבניות השניוניות הריבועיות ומשוואות פל אלמנטריות;

\medskip
\noindent \underline{חלק ב'}: תורת התבניות הדו-צדדיות;

\medskip
\noindent \underline{חלק ג'}: תבניות חלוקת המעגל (כתבניות דו-צדדיות)

\medskip
\noindent \underline{חלק ד'}: תבניות שניוניות מסדר כלשהו, ומשוואות פל מסדר גבוה.

\bigskip
\begin{flushleft}
המחבר
\end{flushleft}

\newpage

% [front matter: table of contents, original page label: IV]
\section*{תוכן העניינים}

\subsection*{פרק א'}
\begin{tabular}{ll}
1--2 & 1. מבוא והגדרות \\
2 & 2. מיוצגי התבנית \\
2--3 & 3. הצגת הבעיה \\
4 & 4. טרנספורמציה ליניארית \\
4--5 & 5. הקשרים בין הבוחנים \\
5--6 & 6. תבניות אקויולנטיות \\
7--9 & 7. סגולות הטרנספורמציה \\
9--10 & 8. תבניות מנוגדות \\
10 & 9. החלוקות התבניות \\
\end{tabular}

\subsection*{פרק ב'}
\begin{tabular}{ll}
11 & 10. הצגות עיקריות \\
11--12 & 11. הבוחן הוא ריבוע מלא \\
12 & 12. המיוצג הוא אפס \\
12--13 & 13. הבוחן אינו ריבוע מלא \\
12--15 & 14. החלוקות ההצגות לפי שיטת גאוס \\
16--19 & 15. חבורות הפתרונות שבקבוצה אחת \\
20--22 & 16. הרחבה למציאת כל הפתרונות השייכים לקבוצה אחת על-פי משוואת פל \\
\end{tabular}

\subsection*{פרק ג'}
% Item 17's text is heavily overtyped/garbled in the source scan.
\begin{tabular}{ll}
23--26 & 17. על משפט אחד [?] הקשור בהעברת ההצגות למחלקה מיוחד[ת] \\
27--32 & 18. קלסיפיקציה חדשה של פתרונות התבניות \\
33--36 & 19. תכונות הפתרונות שבקבוצה אחת, לפי הקלספיקציה החדשה \\
28--32 & 20. צמצום פתירת התבניות לגבי מיצג כלשהו לפתירת תבנית היחידה \\
\end{tabular}

\subsection*{פרק ד'}
\begin{tabular}{ll}
37--45 & 21. פתירת תבנית היחידה, כשהבוחן הוא שלילי \\
\end{tabular}

\subsection*{פרק ה'}
\begin{tabular}{ll}
46 & 22. על קונגרואנציות ממעלה שנייה \\
46--49 & 23. קונגרואנציות לא מלאות \\
50 & 24. פתרון הקונגרואנציות לפי גאוס \\
50--51 & 25. קונגרואנציות דו-צדדיות \\
51--55 & 26. משפטים על הקונגרואנציות הבינומיאליות \\
55--56 & 27. ליקוי בשיטת גאוס \\
\end{tabular}

\newpage
% [original page label: VI]
\subsection*{פרק ו'}
\begin{tabular}{ll}
$[?]$--57 & 28. סקירה כללית על משוואת פל \\
58--61 & 29. הקדמה \\
62 & 30. הוכחה אריתמטית למציאות פתרון משוואת פל \\
63--65 & 31. משפטים יסודיים על משוואת פל \\
66--69 & 32. משפט עזר \\
69 & 33. הקשר של פואנקרה \\
70 & 34. מבוא להכללת משוואת פל \\
73--80 & 35. הכללה א' של משוואת פל; המשוואה מהצורה $ax^2 - by^2 = \pm 1$ \\
80 & 36. הכללת משפט ריכ[?]לדה-פלמשטרים $[?]$; המשוואה מהצורה $2x^2 - Dy^2 = 1$ \\
81 & 38. מקרים מיוחדים \\
83 & 39. הנוסחאות הכלליות למשוואת פל של ההכללה \\
 & 40. הכללה ב' למשוואת פל; המשוואה מהצורה $ax^2 - by^2 = \pm 2$ (עמ' 84) \\
85--86 & 41. מקרים מיוחדים \\
86 & 42. הנוסחאות הכלליות \\
87 & 43. משפט כללי \\
 & 44. הכללה ג' למשוואת פל; המשוואה מהצורה $ax^2 - by^2 = \pm 4$ (עמ' 88) \\
88--$[?]$ & 45. הקשר בין ההכללה הראשונה והשלישית \\
$[?]$ & 46. מקרים מיוחדים \\
90--91 & 47. הנוסחאות הכלליות \\
\end{tabular}

\newpage
% [original page label: VII (partially illegible in scan)]
\subsection*{פרק ז'}
\begin{tabular}{ll}
92--93 & 48. מבוא להכללה הרביעית של משוואת פל - משוואות מסדר גבוה \\
 & 49. על המשוואות מהצורה $ax^{2M} - by^{2N} = \pm 1$ (עמ' 92--93) \\
98--9$[?]$ & (המשך סעיף 49) \\
99--101 & 50. מקרים מיוחדים \\
102--103 & 51. מקרים מיוחדים נוספים \\
 & 52. על המשוואה $x^{2n} - Dy^{2n} = 1$ \\
104--106 & ואי-אפשרות פתרונה בתנאים מיוחדים \\
\end{tabular}

\newpage

% [original page: 1 (unnumbered chapter-opening page)]
\section*{פרק א'}

\subsection*{1. מבוא}

\noindent 1. בשם תבנית ריבועית שניונית אחידה (binäre homogene quadratische Form) אנו קוראים לביטוי מהצורה

\[ ax^2 + bxy + cy^2 \tag{1} \]

במקום שהמקדמים $a$, $b$, $c$ הם מספרים שלמים קבועים ו-$x$, $y$ הם משתנים שלמים.

\medskip
לפי גאוס מסמנים את התבנית דלעיל - עד כמה שאין הדבר נוגע לנעלמים עצמם - ע"י:

\[ (a,\, b,\, c) \tag{2} \]

את הגודל

\[ D = b^2 - 4ac \tag{3} \]

אנו קוראים בשם \underline{בוחן} התבנית (Determinante, Discriminante).

\medskip
לגודל זה בעל התפקידים הרבים באלגברה גם כאן בתורת התבניות השייכת לתורת המספרים חשיבות רבה, כפי שנראה בהמשך התורה הזאת.

\medskip
\noindent\textbf{2. $D$ אינו ריבוע מלא.}

רצינות החקירה על דבר התבניות הריבועיות מתחילה רק כאשר $D$ אינו ריבוע מלא, הואיל ובמקרה אחר התבנית עצמה מתפרקת לשני גורמים ממעלה ראשונה וכל הבעיות שיועלו להלן על הפרק נפתרות במקרה זה בדרך הפשוטה ביותר.

\medskip
לפיכך עלינו לעשות לנוחות הנחה מראש, כי:

\[ \boxed{D \text{ אינו ריבוע מלא}} \tag{4} \]

כמו"כ עלינו להניח מראש, כי:

\[ a \neq 0 \, , \quad c \neq 0 \tag{5} \]

% [original page: 2]
כי במקרה אחר הבעיות תפתרנה טוב - כפי שנראה - בקלות. אולם הנחה כזאת באופן מיוחד היא כבר מיותרת, הואיל היא נובעת מן ההנחה הקודמת. כי לו היה

\[ a = 0 \quad \text{או} \quad c = 0 \]

היה נופק מיד, כי:

\[ D = b^2 \]

בניגוד ל-(4).

\subsection*{3. מיוצגי התבנית}

כל תבנית ריבועית כזאת מייצגת כל מיני מספרים שלמים - לעת עתה המדובר הוא רק בערכים שלמים - ע"י הצבת ערכים שונים במקום $x$, $y$. אם קיים הסווין:

\[ ax^2 + bxy + cy^2 = m \tag{6} \]

בשביל $x$, $y$ מסויימים אנו אומרים שהמספר $m$ מיוצג ע"י התבנית $(a,\,b,\,c)$.

\medskip
כאשר המספרים $x$, $y$ מקיימים את התנאי:

\[ (x,\, y) = 1 \tag{7} \]

(קרי: גדול מחלקיהם המשותף הוא 1, לפי לנדאו)

\medskip
אנו אומרים שההצגה היא \underline{עיקרית} (eigentliche), ובכל מקרה אחר - ההצגה נקראת \underline{לא-עיקרית} (uneigentliche).

\subsection*{4. הצגת הבעייה}

אם כי תורת התבניות על כל אגפיה היא אחת התורות המפותחות ביותר, הרי אפשר בכל זאת לתת סיכום קצר על הבעיות הראשיות אשר מהן נובעות כבר כל שאר הבעיות לכל סוגיהן; יתר על כן: עם פתירת הבעיות הקודמות גם פתרון שאר הבעיות מתקבל כפרי בשל מעל העץ.

\medskip
את הבעיות הראשיות נוכל לפרט כדלקמן:

\medskip
\noindent א) למצוא שיטה שעל פיה אפשר לקבוע ע"י מספר סופי של פעולות אם מספר $m$ מסויים ניתן להצגה ע"י תבנית $(a,\,b,\,c)$

% [original page: 3]
\noindent נתונה.

\medskip
\noindent ב) למצוא $x$, $y$ המקיימים את ההצגה (6) לגבי מספר $m$ מסויים נתון ע"י תבנית מהצורה (1) הנתונה אף היא.

\medskip
\noindent ג) לקבוע את כל המספרים $x$, $y$ המקיימים את ההצגה הנ"ל, במקרה שמספרם הוא סופי.

\medskip
\noindent ד) במקרה שמספר כל המספרים $x$, $y$ הוא אין-סופי - למיינם לקבוצות, ולקבוע את התכונות האפייניות לקבוצות הנ"ל.

\medskip
\noindent ה) לקבוע את התכונות המיוחדות של המייצגים $m$ של תבנית מסויימת עפ"י תכונות שמורות של התבנית, כגון: הבוחן $D$ ומחלקיו.

\medskip
\noindent ו) לקבוע יחסים הדדיים בין התבניות השונות לבין עצמן, באופן בלתי תלוי מן המייצגים $m$, כגון: יחס האקוויולנציה, עפ"י הגדרה מיוחדת שתבוא להלן - וגם באופן תלוי מן המייצג $m$, כפי שזה יבואר עוד.

\medskip
\noindent ז) למיין את התחבניות ע"י חלוקתן למחלקות, מחלקות עפ"י שיטה מסויימת.

\medskip
\noindent ח) לקבוע את מספר המחלקות הנ"ל לגבי $D$ מסויים.

\bigskip
\hrulefill

% [original page: 4]
\subsection*{2. טרנספורמציות ליניאריות}

\noindent 1. תהי

\[ (a,\, b,\, c) \tag{1} \]

תבנית ריבועית נתונה.

\medskip
אם נציב

\[ x = \alpha x' + \beta y' \, , \quad y = \gamma x' + \delta y' \tag{2} \]

אל תוך התבנית הנתונה ונפתחה לפי המשתנים החדשים $x'$, $y'$, אזי נקבל תבנית ריבועית חדשה וגם היא שניונית:

\[ (a',\, b',\, c') \tag{3} \]

ויתקיימו הקשרים הבאים:

\[
\begin{cases}
a' = a\alpha^2 + b\alpha\gamma + c\gamma^2 \\
b' = 2a\alpha\beta + b(\alpha\delta + \gamma\beta) + 2c\gamma\delta \\
c' = a\beta^2 + b\beta\delta + c\delta^2
\end{cases}
\tag{4}
\]

כפי שקל להיווכח.

\medskip
התבנית (3) מתקבלת, איפוא, ע"י טרנספורמציה ליניארית מהתבנית (1). לפי גאוס אנו אומרים, שהתבנית (3) מוכלת בתבנית (1) ולהיפך: התבנית (1) מכילה את התבנית (3) בתוכה.

\medskip
אם נסמן את התבנית (1) ב-$F$ ואת התבנית (3) ב-$F'$ נבטא את העובדה דלעיל ע"י הסימן:

\[ F \supset F' \]

ובעזרת מקדמי ההצגה (2) ע"י:

\[
\begin{pmatrix} \alpha & \beta \\ \gamma & \delta \end{pmatrix} (a,\,b,\,c) = (a',\,b',\,c')
\]

\subsection*{2. הקשרים בין הבוחנים}

ע"י חישובים פשוטים אפשר להיווכח מיד שקיים הקשר הבא, בין הבוחן $D$ של התבנית (1) לבין הבוחן $D'$ של התבנית (3):

% [original page: 5]
\[ D' = (\alpha\delta - \gamma\beta)^2 D \tag{5} \]

ואמנם

\[
D' = b'^2 - 4a'c' = \left[2a\alpha\beta + b(\alpha\delta+\gamma\beta) + 2c\gamma\delta\right]^2 - 4(a\alpha^2+b\alpha\gamma+c\gamma^2)(a\beta^2+b\beta\delta+c\delta^2) =
\]
\[
= (\alpha\delta-\gamma\beta)^2(b^2-4ac) = (\alpha\delta-\gamma\beta)^2 D .
\]

מכאן נמצינו למדים, שלא כל שתי תבניות ניתנות להעברה אחת לשניה ע"י טרנספורמציה ליניארית.

\medskip
\underline{התנאי ההכרחי} לכך, שמבין שתי תבניות, האחת תכיל את השניה הוא קיום הקשר:

\[ D' = t^2 D \tag{6} \]

בין הדיסקרימיננטות $D$, $D'$ של התבניות הללו, במקום ש-$t$ הנו שלם.

\medskip
מאידך, אין תנאי זה מספיק. ואמנם, יתכן, שמתקיים הקשר \underline{(6)} בין הדיסקרימיננטות של התבניות ובכל זאת אין האחת מהן מכילה את השנייה. יתר על כן: יתכן שקיים שוויון בין הדיסקרימיננטות:

\[ D' = D. \tag{7} \]

ובכל זאת אין האחת ניתנת להעברה לשניה ע"י טרנספורמציה ליניארית מסוג (2), כפי שהדבר יתברר עוד.

\subsection*{3. תבניות אקוויולנטיות}

כפי שראינו כבר מתבטא הקשר שבין הדיסקרימיננטות של שתי תבניות, שהאחת מוכלת בשניה בעזרת דטרמיננט ההצגה:

\[
\begin{vmatrix} \alpha & \beta \\ \gamma & \delta \end{vmatrix} = \alpha\delta - \gamma\beta = t \tag{8}
\]

במקום ש-$t$ הנו מספר שלם מסויים.

\medskip
כאשר $t$ הנו מספר חיובי אנו אומרים לפי גאוס, שהתבנית (1) מכילה את התבנית (3) באופן \underline{עיקרי}, וכאשר $t$ הנו מספר שלילי אנו אומרים שהתבנית (1) מכילה את התבנית (3) באופן \underline{לא-עיקרי}.

% [original page: 6]
מן ההצגה (2) אנו מקבלים בלי קושי את הטרנספורמציה ההפוכה, והיא:

\[
\begin{cases}
x' = \dfrac{1}{\alpha\delta - \gamma\beta} (\delta x - \beta y) \\[2mm]
y' = \dfrac{1}{\alpha\delta - \gamma\beta} (-\gamma x + \alpha y)
\end{cases}
\tag{9}
\]

מכאן אנו רואים, שאין הטרנספורמציה ההפוכה מתקבלת במקדמים שלמים, כאשר $|t| > 1$, לפיכך, אם $F \supset F'$ אין $F' \supset F$ כאשר $|t| > 1$.

\medskip
לעומת זאת יתכן מקרה שבו יתקיימו שני התנאים בבת-אחת: $F \supset F'$, $F' \supset F$.

\medskip
את מקרה זה אנו מבטאים בסימן: $F \sim F'$

\medskip
ויש מקום למקרה זה אך ורק כאשר:

\[ t = \pm 1 \, ; \tag{10} \]

כי הטרנספורמציה ההפוכה מקבלת אז את הצורה:

\[ x' = \pm(\delta x - \beta y) \, , \quad y' = \pm(-\gamma x + \alpha y) \tag{11} \]

וע"כ אפשר להעביר בחזרה את התבנית (3) ע"י הטרנספורמציה (11) בעזרת מקדמים שלמים.

\medskip
אנו אומרים במקרה זה לפי גאוס, שהתבניות הן אקוויולנטיות זל"ז, ומשתמשים בסימן:

\[
\begin{cases} (a,\,b,\,c) \sim (a',\,b',\,c') \\ \text{כאשר } t = 1 \end{cases}
\tag{12}
\]

ובסימן:

\[
\begin{cases} (a,\,b,\,c) \approx (a',\,b',\,c') \\ \text{כאשר } t = -1 \end{cases}
\tag{13}
\]

במקרה ראשון נגיד, שהתבניות הן אקוויולנטיות באופן \underline{עיקרי}, ובמקרה שני, שהן אקוויולנטיות באופן \underline{לא-עיקרי}.

% [original page: 7]
\subsection*{4. כללים יסודיים}

כאשר $F \supset F'$ ו-$F' \supset F''$ ניווכח בנקל, שגם $F \supset F''$; ובכן: היחס $\supset$ הוא טרנסיטיבי; ואמנם מן ההצגות:

\[ x = \alpha x' + \beta y' \, , \quad y = \gamma x' + \delta y' \, , \quad \alpha\delta - \gamma\beta = t \tag{14} \]
\[ x' = \alpha' x'' + \beta' y'' \, , \quad y' = \gamma' x'' + \delta' y'' \, , \quad \alpha'\delta' - \gamma'\beta' = t' \tag{15} \]

נופק מיד, כי:

\[
\begin{cases}
x = (\alpha\alpha'+\beta\gamma')x'' + (\alpha\beta'+\beta\delta')y'' \\
y = (\gamma\alpha'+\delta\gamma')x'' + (\gamma\beta'+\delta\delta')y''
\end{cases}
\tag{16}
\]

וחוץ מזה מתקיים:

\[
(\alpha\alpha'+\beta\gamma')(\gamma\beta'+\delta\delta') - (\alpha\beta'+\beta\delta')(\gamma\alpha'+\delta\gamma') = (\alpha\delta-\gamma\beta)(\alpha'\delta'-\gamma'\beta') = tt' . \tag{17}
\]

לפיכך, אם $(a,\,b,\,c) \supset (a',\,b',\,c')$ ו- $(a',\,b',\,c') \supset (a'',\,b'',\,c'')$ אזי:

\[ (a,\,b,\,c) \supset (a'',\,b'',\,c'') \]

וחוץ מזה קיים:

\[
\begin{cases}
D' = t^2 D. \\
D'' = b''^2 - 4a''c'' = t'^2 D' = (tt')^2 D
\end{cases}
\tag{18}
\]

אולם כאשר $t = \pm 1$

% [original page: 8]
נופק מיד, כי:

\[ D'' = D' = D. \tag{19} \]

לגבי תבניות אקוויולנטיות קל להיווכח, שקיימים הכללים הבאים:

\medskip
\noindent א) $(a,\,b,\,c) \sim (a,\,b,\,c)$

ואמנם $\begin{pmatrix}1 & 0\\ 0 & 1\end{pmatrix}(a,\,b,\,c) = (a,\,b,\,c)$ כפי שנובע מחשבונות קלים;

\medskip
\noindent ב) אם $(a,\,b,\,c) \sim (a',\,b',\,c')$ אזי: $(a',\,b',\,c') \sim (a,\,b,\,c)$

ובכן: \underline{יחס האקוויולנציה הוא יחס סימטרי}.

\medskip
ואמנם, אם $\begin{pmatrix}\alpha & \beta\\ \gamma & \delta\end{pmatrix}(a,\,b,\,c) = (a',\,b',\,c')$ אזי: $\begin{pmatrix}\delta & -\beta\\ -\gamma & \alpha\end{pmatrix}(a',\,b',\,c') = (a,\,b,\,c)$

כפי שנובע מנוסחה (9).

\medskip
וכן ב'): אם $(a,\,b,\,c) \approx (a',\,b',\,c')$ אזי $(a',\,b',\,c') \approx (a,\,b,\,c)$ כי אם $\alpha\delta - \gamma\beta = -1$ אזי, גם: $\delta\alpha - (-\gamma)(-\beta) = -1$

\medskip
\noindent\textbf{\underline{הערה.}} יתכן גם ששתי תבניות תהיינה בבת-אחת אקוויולנטיות באופן עיקרי ובאופן לא-עיקרי. כאלה הן, למשל, התבניות: $(1,\,0,\,-2)$, $(-1,\,0,\,2)$

% [original page: 9]
כפי שאפשר להיווכח מן ההצגות בעלות הדטרמיננטות

\[
\begin{vmatrix} 1 & 2 \\ -1 & -1 \end{vmatrix}
\qquad
\begin{vmatrix} 1 & 2 \\ 1 & 1 \end{vmatrix}
\]

ואמנם:

\[
\begin{pmatrix} 1 & 2 \\ -1 & -1 \end{pmatrix}(1,\,0,\,-2) = (-1,\,0,\,-2)
\]
\[
1\cdot(-1) - (-1)\cdot 2 = 1
\]

וכן:

\[
\begin{pmatrix} 1 & 2 \\ 1 & 1 \end{pmatrix}(1,\,0,\,-2) = (-1,\,0,\,-2)
\]
\[
1\cdot 1 - 2\cdot 1 = -1 .
\]

\medskip
\noindent ג) אם $(a,\,b,\,c) \sim (a'',\,b'',\,c'')$ ו- $(a',\,b',\,c') \sim (a'',\,b'',\,c'')$ אזי: $(a,\,b,\,c) \sim (a',\,b',\,c')$

\medskip
\noindent ד) אם $(a,\,b,\,c) \approx (a'',\,b'',\,c'')$ ו- $(a',\,b',\,c') \approx (a'',\,b'',\,c'')$ אזי: $(a,\,b,\,c) \sim (a',\,b',\,c')$.

\medskip
בנכונות הכללים האחרונים אפשר להיווכח באופן יותר יעיל, על סמך הנוסחאות (14), (15), (16), (17).

\subsection*{5. תבניות מנוגדות (Entgegengesetzte Formen)}

שתי תבניות מהצורה $(a,\,b,\,c)$, $(a,\,-b,\,c)$ נקראות, לפי גאוס, בשם \underline{תבניות מנוגדות}.

\medskip
גאוס קבע, ששתי תבניות כאלו הן אקוויולנטיות באופן לא-עיקרי. ואמנם ע"י הטרנספורמציה $\begin{pmatrix}1 & 0\\ 0 & -1\end{pmatrix}$ עוברת התבנית $(a,\,b,\,c)$ לתבנית

% [original page: 10]
$(a,\,-b,\,c)$, ואכן קיים: $1\cdot(-1) - 0\cdot 0 = -1$

\subsection*{6. התחלקות התבניות}

כפי שראינו, התבניות האקוויולנטיות הן בנות דיסקרימיננטות שוות; אולם, כפי שכבר העירונו בסעיף הקודם, אין כל התבניות אשר להן דיסקרימיננטות שוות, אקוויולנטיות - לא באופן עיקרי ולא באופן לא-עיקרי.

\medskip
לפיכך נוהגים לחלק את כל התבניות בנות דיסקרימיננטה $D$ אחת למחלקות, מחלקות של תבניות אקוויולנטיות זל"ז.

\medskip
מאחר שע"י טרנספורמציה בעלת דטרמיננטה

\[
\begin{vmatrix} \alpha & \beta \\ \gamma & \delta \end{vmatrix} \, , \quad \alpha\delta - \gamma\beta = 1
\]

אפשר לקבל תבניות אקוויולנטיות במספר אין-סופי ברור, שכל מחלקה כזאת מכילה הרבה עד אין-סוף תבניות.

\medskip
מהו מספר המחלקות הללו?

\medskip
זאת היא אחת הבעיות הקשות ביותר של תפרת התבניות ולא נשיב לשאלה זאת כאן.

\medskip
השאלה השנייה במעלה היא: איך לקבוע אם ומתי שתי תבניות הן אקוויולנטיות זל"ז?

\medskip
לשאלה זאת נקדיש את אחד הפרקים להלן.

\bigskip
\hrulefill

% [original page: 11]
\section*{פרק ב'}

\subsection*{1. התחלקות ההצגות $(x,\,y)$ לגבי $m$ מסויים (לפי שיטת גאוס)}

\noindent 1. יהי:

\[ ax^2 + bxy + cy^2 = m \tag{1} \]

מותר להניח, כי:

\[ (x,\,y) = 1 \tag{2} \]

הואיל ובמקרה

\[ (x,\,y) = d \neq 1 \]

קיים

\[ x = dx' \, , \quad y = dy' \]

בשביל $x'$, $y'$ שלמים מסויימים, ולכן:

\[
d^2(ax'^2 + bx'y' + cy'^2) = m
\]

ומכאן, כי:

\[ d^2 \mid m \]

כלומר

\[ m = d^2 m' \]

והצווין (1) מצטמצם לצורה:

\[ ax'^2 + bx'y' + cy'^2 = m' \]

שבה מתקיים כבר התנאי:

\[ (x',\,y') = 1. \]

כן מותר להניח, כי:

\[ (a,\,b,\,c) = 1 \tag{3} \]

(קרי: גדול מחלקיהם המשותפים של $a$, $b$, $c$ הוא 1), מהסיבה הקודמת.

\subsection*{2. $D$ ריבוע מלא}

במקרה זה מתפרקת, כידוע, התבנית לשני גורמים ליניאריים:

\[ (ex+fy)(e'x+f'y) \tag{4} \]

% === TRANSCRIPTION NOTES ===
% - p000 (title page): year read as "תש\"ל" (5730 = 1970) after close
%   zoom inspection; an earlier, blurrier pass misread it as "תש\"י".
%   Library stamps and handwritten accession marks ("416308", "2X")
%   present on the page were omitted from the body per instructions.
% - p001 (dedication): handwritten marks, a signature, and the date
%   "6.7.50" at the top of the page were not transcribed (handwritten
%   marginalia). One library "withdrawn from stock" stamp also present,
%   omitted.
% - p002 (introduction, page I): one short phrase in the paragraph
%   beginning "וזאת היא אחת הסיבות..." is badly overtyped/smudged in
%   the scan; transcribed with [?] markers as best-effort
%   ("ה[?] ראשונים [?]"). Likewise the phrase "בפוריות ה[?]ישרית [?]
%   הרידונקלית" later on the page is uncertain and marked with [?].
%   The name "דיריכלה" (Dirichlet) appears letter-spaced in the
%   original (typewriter emphasis), rendered here in \textbf{}.
% - p003: the name "גאוס" (Gauss) is printed in the source as "גאום"
%   (with a final mem rather than samekh); silently normalized to the
%   standard spelling "גאוס" throughout this chunk. The term
%   "עדדיות"/"צדדיות" fluctuates in the scan; standardized throughout
%   to "צדדיות" ("דו-צדדיות") as the evident intended reading.
%   Occurrences of "שכיוניות" were normalized to "שניוניות" (secondary/
%   quadratic) to match the book's title and consistent usage elsewhere.
% - p005-p008 (table of contents): several page-range numbers and a
%   few item titles are faint or overtyped in the scan; uncertain
%   digits/words are marked with [?] inline (see chapter ו' and ז'
%   entries, e.g. the start page of item 28, the second digit of the
%   "98-9[?]" range, item 36's title, and the page numbers for items
%   45-46). The original page-label Roman numerals for pages
%   corresponding to p003, p006, and p008 were faint/uncaptured in the
%   scan; II, V, and VII were inferred from the surrounding sequence
%   (I, III, IV, VI visible).
% - p009-p019 (Chapter A, sections 1-2, opening of Chapter B): text and
%   formulas are clearly legible typewritten/handwritten-symbol content;
%   no illegible passages encountered. Greek letters alpha/beta/gamma/
%   delta were hand-inserted in the original scan and are rendered here
%   as \alpha, \beta, \gamma, \delta.
